% ============================================================
% APPENDIX F — Autonomous Final Project: BookShelf
% ============================================================
\chapter{Autonomous Final Project: BookShelf}
\label{app:bookshelf}

\section*{Your Architecture Exam}

This project has no step-by-step tutorial. You have the brief, the starter \texttt{\_CONTEXT.md}, and the professional workflow from Section~3.8. Your task is to apply the method from start to finish: generate specifications, derive use cases, structure epics, break them down into tasks, and implement.

If you can complete BookShelf on your own, you have internalized the ADLC method.

% ---------------------------------------------------------
\section{The Brief}

\textbf{BookShelf} is an application for managing a personal library. You can catalog books you own, books you are reading, and books you want to read. The application includes an API backend, a web frontend, and a mobile app.

\subsection{Functional Requirements}
\begin{enumerate}
  \item \textbf{Book catalog}: full CRUD (title, author, year, genre, optional ISBN, optional cover URL)
  \item \textbf{Reading status}: each book can be ``To Read'', ``Reading'', ``Read'', or ``Abandoned''
  \item \textbf{Personal notes}: users can add notes and a rating from 1 to 5 stars
  \item \textbf{Search and filters}: by title/author, status, and genre
  \item \textbf{Statistics}: counts by status, books read per year, and most-read genre
  \item \textbf{Authentication}: Google OAuth~2.0 login
  \item \textbf{Multi-user}: each user must see only their own library
\end{enumerate}

\subsection{Non-Functional Requirements}
\begin{itemize}
  \item Backend: Node.js + Express + PostgreSQL + Prisma
  \item Frontend: React + Tailwind CSS
  \item Mobile: Flutter
  \item Deployment: Railway, Vercel, Neon, or equivalent platforms
\end{itemize}

% ---------------------------------------------------------
\section{The Starter \texttt{\_CONTEXT.md}}

This is your starting point. Copy it to the project root, then improve it as you learn.

\begin{lstlisting}[language={},title={\_CONTEXT.md --- BookShelf starter}]
# Project: BookShelf

## Purpose
Web + mobile application for managing a personal library.
Each user can catalog books, track reading status,
add notes, and view statistics.

## Technologies
- Backend: Node.js 20+, Express.js, PostgreSQL 16, Prisma ORM
- Frontend: React 18+ with Vite, Tailwind CSS
- Mobile: Flutter 3.x with Dart, Riverpod for state management
- Auth: OAuth 2.0 with Google (passport.js)
- Tokens: JWT (access 1h, refresh 7d)

## Structure
bookshelf/
├── _CONTEXT.md
├── backend/
│   ├── src/
│   │   ├── routes/
│   │   ├── controllers/
│   │   ├── services/
│   │   ├── middleware/
│   │   └── utils/
│   ├── prisma/
│   │   └── schema.prisma
│   └── tests/
├── frontend/
│   ├── src/
│   │   ├── components/
│   │   ├── pages/
│   │   ├── hooks/
│   │   ├── services/
│   │   └── context/
│   └── tests/
└── mobile/
    └── bookshelf_app/
        └── lib/
            ├── models/
            ├── providers/
            ├── screens/
            ├── services/
            └── widgets/

## Data Model (indicative Prisma schema)
- User: id, email, name, avatar, createdAt
- Book: id, userId, title, author, year, genre, isbn?, coverUrl?,
        status (TO_READ | READING | READ | ABANDONED),
        rating (1-5)?, notes?, createdAt, updatedAt
- ReadingSession: id, bookId, startDate, endDate?, pagesRead?

## Conventions
- Naming: camelCase for JS/TS, snake_case for DB columns
- API responses: { success: boolean, data: T, error?: string }
- Validation: Zod on all endpoints
- Commits: Conventional Commits

## Constraints (SEC/PERF)
- SEC-01: JWT secret >= 256 bits, stored in environment variable
- SEC-02: Every query filters by userId — NO exceptions
- SEC-03: Input sanitized and validated before reaching the database
- SEC-04: Rate limiting on auth endpoints (10 req/min)
- SEC-05: Helmet.js for security headers
- PERF-01: Pagination on GET /books (default 20, max 100)
- PERF-02: Index on Book.userId and Book.status
\end{lstlisting}

% ---------------------------------------------------------
\section{The Workflow: From Idea to Code}

This project has no guided solution, but it does have a \textbf{method}. Apply the professional workflow from Section~3.8 to break the work into manageable units.

\subsection{Step 1: Generate the specifications}
\begin{lstlisting}[language={}]
"Here is the BookShelf brief: [paste F.1].
 Generate a complete PRD with functional requirements (FR-01, FR-02…),
 non-functional requirements (NFR-01…), technical constraints, and priorities.
 Save to docs/PRD.md."
\end{lstlisting}

\subsection{Step 2: Derive the use cases}
\begin{lstlisting}[language={}]
"Based on the PRD in docs/PRD.md, generate the main use cases.
 For each use case: actor, precondition, main flow,
 alternative flows, postcondition. Save to docs/USE_CASES.md."
\end{lstlisting}

\subsection{Step 3: Design and technical decisions}
\begin{lstlisting}[language={}]
"Based on the PRD and use cases for BookShelf, propose the system
 architecture. For each technical decision (database, ORM, authentication,
 project structure, API patterns) present 2-3 options with pros, cons,
 and a recommendation. I choose, you document in docs/DESIGN.md."
\end{lstlisting}

\subsection{Step 4: Structure the epics}
\begin{lstlisting}[language={}]
"Based on the use cases, organize the work into epics (E-001, E-002…).
 For each epic: title, objective, covered use cases, dependencies.
 Save the index to docs/epics/INDEX.md and one file per epic."
\end{lstlisting}

\subsection{Step 5: Break the work into tasks}
\begin{lstlisting}[language={}]
"Break down epic E-003 into atomic tasks (T-003.1, T-003.2…).
 Each task must have: verifiable acceptance criteria,
 involved files, applicable SEC/PERF constraints, story points."
\end{lstlisting}

\subsection{Step 6: Implement task by task}
\begin{lstlisting}[language={}]
"Implement T-003.1. Read docs/epics/tasks/E-003/T-003.1_*.md.
 Follow the constraints from _CONTEXT.md."
\end{lstlisting}

\begin{tipbox}
Do not move to the frontend before the backend tasks are complete and verified. The task breakdown gives you visibility --- use it.
\end{tipbox}

% ---------------------------------------------------------
\section{Phase Challenges}

\subsection{Phase 1: Backend (Chapters 6--8)}
\begin{enumerate}
  \item Initialize the Node.js project with Express
  \item Configure Prisma with the \texttt{Book} and \texttt{User} schemas
  \item Implement CRUD endpoints for books
  \item Add Google OAuth~2.0 authentication
  \item Add the statistics endpoint (\texttt{GET /api/stats})
  \item Write unit and integration tests
  \item Verify the security checklist (SEC-01 $\rightarrow$ SEC-05)
\end{enumerate}

\subsection{Phase 2: Frontend (Chapters 9--10)}
\begin{enumerate}
  \item Create the React project with Vite and Tailwind
  \item Implement the login page with OAuth redirect
  \item Build the dashboard with book list and filters
  \item Implement the form to add or edit books
  \item Create the statistics page with counts and simple charts
  \item Integrate with the real backend
  \item Write tests with Testing Library
\end{enumerate}

\subsection{Phase 3: Mobile (Chapters 11--12)}
\begin{enumerate}
  \item Create the Flutter project
  \item Implement mobile authentication
  \item Build the book list with pull-to-refresh
  \item Add the book detail page with notes and rating
  \item Implement search and filters
  \item Write widget tests
\end{enumerate}

\subsection{Phase 4: Production (Chapters 14--15)}
\begin{enumerate}
  \item Run an OWASP Top~10 audit on the backend
  \item Apply the security fixes
  \item Deploy the backend to Railway
  \item Deploy the frontend to Vercel
  \item Host the database on Neon
  \item Verify end-to-end behavior
\end{enumerate}

% ---------------------------------------------------------
\section{Success Criteria}

\begin{center}
\begin{tabularx}{\textwidth}{c X}
\toprule
\textbf{Status} & \textbf{Criterion} \\
\midrule
$\square$ & A user can log in with Google \\
$\square$ & They can add, edit, and delete books \\
$\square$ & They can change the reading status \\
$\square$ & They can add notes and ratings \\
$\square$ & They can search and filter the catalog \\
$\square$ & They can view statistics on the dashboard \\
$\square$ & The mobile app shows the same data \\
$\square$ & All tests pass \\
$\square$ & The security checklist is completed \\
$\square$ & The application is accessible online \\
$\square$ & The \texttt{\_CONTEXT.md} is updated with lessons learned \\
\bottomrule
\end{tabularx}
\end{center}

% ---------------------------------------------------------
\section{Hints}

\begin{itemize}
  \item Follow the workflow from Section~3.8: PRD $\rightarrow$ Use Cases $\rightarrow$ Epics $\rightarrow$ Tasks $\rightarrow$ Implementation.
  \item Start from the backend. It is the foundation for everything else.
  \item Use the \texttt{\_CONTEXT.md} from the very beginning.
  \item Test each phase before moving on to the next one.
  \item If the AI generates something you do not understand, ask it to explain the code line by line.
  \item If you get stuck, go back to the corresponding chapter of the book and reuse the same method.
\end{itemize}

\vfill
\begin{center}
\textit{If you complete BookShelf on your own, congratulations: you are no longer just using AI tools --- you are working as a context architect.}
\end{center}
